\chapter*{Abstract}

\addcontentsline{toc}{chapter}{Abstract}

Railway alignment modelling remains a central task in infrastructure
engineering, yet existing approaches are often fragmented across
geometric design, data standards, and dynamic evaluation.

This work introduces the concept of
\emph{Alignment-Based Information Modelling (AIM)},
a unified framework in which railway alignment is represented as a
state-based object combining geometry, curvature, and cant.

Within this framework, alignment is not treated as a static curve, but
as a dynamic system governed by curvature functions and their
derivatives.
Transition curves are reinterpreted as control functions for curvature,
rather than as predefined geometric templates.

A canonical representation is proposed, integrating horizontal
geometry, vertical behaviour, and cant into a single model.
This representation enables direct evaluation of dynamic quantities
such as effective acceleration and jerk, and provides a natural basis
for numerical reconstruction and optimization.

The work further establishes a connection between alignment modelling
and constrained optimization.
Classical transition curves, such as the clothoid, are derived as
solutions of variational problems in curvature space.
More general formulations allow the inclusion of engineering
constraints and objective functions reflecting comfort, safety, and
cost considerations.

To bridge the gap between computational modelling and real-world data,
the concept of a normalization pipeline is introduced.
Heterogeneous input formats, including LandXML, IFC, and legacy data,
are transformed into a canonical alignment representation.
An interactive intermediate layer supports user-guided interpretation
of imported data.

Building on this foundation, the notion of
\emph{BIMalignment} is proposed as an extension of BIM principles.
In this view, alignment becomes the central organizing structure for
infrastructure models, linking geometry, semantics, and dynamic
behaviour within a computationally enabled framework.

The presented approach is demonstrated through numerical examples,
optimization sketches, and application-oriented use cases.
The results indicate that alignment-centred modelling provides a
coherent and extensible basis for future railway design systems.

The work concludes by outlining directions for further research,
including advanced vehicle dynamics, network-level optimization,
real-time interaction, and integration with AI-assisted design tools.

Overall, the AIM framework establishes a shift from static geometric
representation toward dynamic, optimization-ready, and
alignment-centred infrastructure modelling.
